\documentclass[11pt,a4paper]{article}
\usepackage[T1]{fontenc}
\usepackage[utf8]{inputenc}
\usepackage{lmodern}
\usepackage[margin=25mm,headheight=15pt]{geometry}
\usepackage{amsmath,amssymb,amsthm,mathtools}
\usepackage{booktabs,array,longtable}
\usepackage{microtype,xcolor,enumitem,fancyhdr}
\usepackage{hyperref}
\hypersetup{colorlinks=true,linkcolor=blue!45!black,citecolor=blue!45!black,urlcolor=blue!55!black,
 pdftitle={Smaller explicit Diophantine equations for ZFC consistency},
 pdfauthor={Research synthesis prepared for Vladimir Reshetnikov}}
\newcommand{\N}{\mathbb N}
\newcommand{\Z}{\mathbb Z}
\newcommand{\Q}{\mathbb Q}
\newcommand{\Zpos}{\mathbb Z_{>0}}
\newcommand{\Con}{\operatorname{Con}}
\newcommand{\pc}{\operatorname{pc}}
\newcommand{\pin}{\texttt{124c07c0b77a9e20e84676514d8b1ff78536ea3f}}
\newtheorem{theorem}{Theorem}[section]
\newtheorem{lemma}[theorem]{Lemma}
\newtheorem{proposition}[theorem]{Proposition}
\newtheorem{corollary}[theorem]{Corollary}
\theoremstyle{definition}
\newtheorem{definition}[theorem]{Definition}
\theoremstyle{remark}
\newtheorem{remark}[theorem]{Remark}
\numberwithin{equation}{section}
\setlist{nosep,leftmargin=*}
\setlength{\parskip}{3pt}
\setlength{\parindent}{0pt}
\setlength{\emergencystretch}{2em}
\pagestyle{fancy}
\fancyhf{}
\fancyhead[L]{\small Explicit Diophantine consistency equations}
\fancyhead[R]{\small ProveIt / 23 October 2026}
\fancyfoot[C]{\thepage}

\begin{document}
\hypersetup{pageanchor=false}
\begin{titlepage}
\vspace*{18mm}
{\small\sffamily\color{blue!45!black} RESEARCH ARTICLE AND REPRODUCIBILITY ARCHIVE}\par
\vspace{9mm}
{\Huge\bfseries Smaller explicit\par Diophantine equations\par for ZFC consistency}\par
\vspace{9mm}
{\Large Encoding only the monomials that occur,\par and recovering a deleted Pell quotient}\par
\vspace{9mm}
{\large Prepared for Vladimir Reshetnikov and the ProveIt project\par
23 October 2026}
\vspace{10mm}
\begin{abstract}
We construct explicit integer-coefficient polynomials whose absence of
zeros, on precisely specified domains, is equivalent to syntactic
consistency of ZFC through the repository's calibrated blank-K215 compiler.
The comparison includes the eight wave-6 reports at the pinned revision,
not just the older collection summary. A support-dependent quadratic
encoding lowers the nine-positive-variable degree from $129\,875\,980$
to $30\,947\,584$, with 651 binary arithmetic operations instead of 788.
Its ten-unrestricted-integer companion has degree $56\,593\,864$ and
653 operations, versus $237\,520\,456$ and 782. The general ten-integer
reduction is Sun's; the contribution is the smaller explicit source
encoding and its specialization.

A separate construction uses eleven positive variables, degree 954,
and 115 operations. It combines a refined auxiliary-root estimate with
a midpoint carrier and the modulus $4q^4$, which recovers the deleted
index quotient without a translation of the scale variable. A lower-degree
fourteen-positive-variable alternative retains degree 56 with 79 operations,
removing one variable from the historical degree-56 construction.
All constants are fixed
by finite recipes and binary payloads; no consistency-dependent choice
or unspecified enumeration index occurs. The archive contains the
literal circuits, complete new reduction proofs, exact finite certificates,
and preserved source notices. The inherited compiler calibration and
number-theoretic interfaces are identified explicitly. This is an ordinary
mathematical construction, not a new end-to-end Lean certificate or a
proof of global minimality.
\end{abstract}
\vfill
{\small Reviewed repository revision: \pin\par
\url{https://github.com/VladimirReshetnikov/ProveIt}}
\end{titlepage}
\pagenumbering{roman}
\hypersetup{pageanchor=true}
\tableofcontents
\clearpage
\pagenumbering{arabic}

\section{Results, comparison, and the exact logical assertion}

\subsection{What has been constructed}

Fix the ordinary syntactic consistency statement
\[
 \Con(\mathrm{ZFC})\;\equiv\;
 \text{there is no finite derivation of falsity from the chosen ZFC axiom predicate}.
\]
The repository identifies this predicate with
$\neg\,\texttt{\detokenize{BProv ZFCax_s [] Form.fBot}}$ at its formal
endpoint. The constructions below are bound to the same literal
215-state machine and one complete native coefficient tuple.

\begin{theorem}[Consistency equations, with stated inherited interfaces]
\label{thm:main}
Use the imported mathematical calibration of the literal blank-K215
compiler stated in Section~\ref{sec:dependencies}, and the cited arithmetic
interfaces, as theorems with ordinary proofs formalizable in ZFC.
For every polynomial $P$ constructed here, on its stated domain $D$,
composition with the displayed reductions yields
\[
 \mathrm{ZFC}\ \vdash\
 \left(\neg\exists\boldsymbol{x}\in D\;P(\boldsymbol{x})=0\right)
 \longleftrightarrow\Con(\mathrm{ZFC}).
\]
The eleven- and fourteen-positive-variable constructions use the
ordinary Pell/compiler interface and do not use Sun's 2026 preprint.
The nine-variable construction uses the repository's Jones interface;
the ten-unrestricted-integer construction additionally uses Sun's
ten-integer interface.
\end{theorem}

This is an ordinary provability claim obtained by composing the stated
proof dependencies; assuming only their external truth would not suffice
for the displayed internal conclusion. No consistency assumption is used.
The new implications are proved below; the full unbounded compiler is
inherited from the pinned repository, with its ordinary proofs included.
We have not produced a new proof-assistant derivation of that compiler
or an object-language ZFC proof term. Section~\ref{sec:formal-status}
keeps these levels distinct.

The four principal outputs improve different measures. They cannot be
ranked by one unqualified use of ``simplest.'' In particular, the shortest
known circuit in the reviewed family remains the historical
67-operation $P_{67}$, with fifteen positive variables and degree 155.

\begin{center}
\small
\begin{tabular}{@{}p{34mm}rrp{48mm}@{}}
\toprule
Domain / unknowns & Exact degree & Operations & Reviewed comparison\\
\midrule
$9$ positive & $30\,947\,584$ & $651$ & Package 45: $129\,875\,980$, 788\\
$9$ nonnegative & $30\,947\,584$ & $654$ & Package 45: same degree, 791\\
$10$ unrestricted & $56\,593\,864$ & $653$ & Package 40: $237\,520\,456$, 782\\
$11$ positive & $954$ & $115$ & Package 38: $1140$, 129; or $1154$, 123\\
$11$ positive, system & $2$; $952$ & $110$ & Two equations, shifted; equality tests free\\
$14$ positive & $56$ & $79$ & Historical $Q_{56}$: $15$ variables, degree 56, 74 operations\\
\bottomrule
\end{tabular}
\end{center}

The fourteen-variable lower-degree alternative and its distinct
tradeoff are given in Section~\ref{sec:fourteen}. The archive also retains
an eleven-variable full-carrier variant of degree 958 at 116 operations,
which provides a simpler comparison with package 44's twelve-variable
carrier. The midpoint variant dominates it on the two displayed measures.

The new nine-variable degree is about $4.20$ times smaller than the
wave-6 comparison, and its evaluator removes 137 operations. The ten-integer
degree is also about $4.20$ times smaller and removes 129 operations
from the degree-best comparison. Package 38 has a different ten-integer
point at 773 operations and larger degree; the present 653-operation
point improves that evaluator as well. The eleven-variable point removes
186 degrees and 14 operations from the previously lowest-degree
eleven-variable equation, or 200 degrees and 8 operations from the
123-operation alternative.

\subsection{Counting conventions and what is new}

An operation is one binary addition, subtraction, or multiplication,
including a multiplication by an integer numeral or a squaring. Shared
registers are evaluated once. A fixed coefficient is a leaf, and its
finite construction is accounted for separately. Fixed powers in a
display are expanded into multiplication gates in the literal circuit.
Thus these are circuit sizes, not counts of monomials in a fully expanded
polynomial. All supplied variables, fixed ports, and emitted gates are live.
For a system, comparisons are free under the same convention used by the
repository. A natural variable includes zero; a positive variable does not.

The concrete contributions relative to the reviewed material are:
\begin{enumerate}
\item Only the 71 quadratic monomials that actually occur are isolated
against all 528 possible quadratic expansion monomials. Full pair-sum
uniqueness is unnecessary. This permits maximum weight 190 instead of 799.
\item Auxiliary-root recovery needs
$d(Ig-cf)>b+cf$, rather than the stronger $Ig-cf>1$.
This removes a circular dependence on the index quotient being recovered.
\item The even polynomial modulus $U=4q^4$ replaces a much larger carrier
scale. Its gcd with $H$ is at most three, and its congruence information
is enough to recover an exact Pell-index equality.
\item The midpoint $\ell+k/2$ can be used in a carrier with integer
coefficients because $U/2=2q^4$ is itself a polynomial. No additional
variable or global doubling is needed.
\item An early index lower bound makes the same small modulus compatible
with the ordinary strong-rank equation. This removes one variable while
retaining degree 56, at the cost of five operations.
\end{enumerate}
Sparse polynomial encodings, Pell equations, and relation-combining
norms are established methods~\cite{jones1982,sun2021,sun2026}.
We claim the stated new combinations,
finite weight assignment, explicit circuits, and comparisons with this
repository revision. We do not claim the first use of support-sensitive
Kronecker encoding in all of mathematics, an optimal weight assignment,
or a global lower bound on the number of variables.

\subsection{The reviewed baseline}

The pinned commit is \pin~\cite{repo}, whose intake commit places wave-6 sources
38--45 beside sources 31--37. At that revision, the collection README
still primarily summarizes wave 5. We therefore inspected the actual
wave-6 source files and literal circuit data. The original incoming
archives have been placed into the report directory; the pinned
\texttt{docs/incoming} directory contains its placement ledger.
The relevant current path is
\begin{center}
\small\texttt{Computability/DiophantineConsistency/con-zfc-diophantine-forms/}.
\end{center}
Historical source paths in preserved files refer to the earlier
\texttt{SetTheory/Cardinals} location at their own named commits. They are
provenance, not instructions to mix old coefficient tuples.

\section{One fixed coefficient tuple and the inherited interfaces}
\label{sec:dependencies}

\subsection{The native coefficients are fixed integers}

Every new main construction uses the unchanged, complete package-31
tuple, whose base is
\[
 \rho=2^{121},\qquad L_c=8801,\qquad
 B_c=\rho^{L_c}=2^{1\,064\,921}.
\]
The four ports are
\[
 \beta=B_c-1,\qquad K_0=D_C+B_cD_R,\qquad M_C,
 \qquad M_F=M_{F,\mathrm{native}}+B_c-1.
\]
The last port is already shifted. Substituting the unshifted high mask
would define a different and unjustified equation.

The exact finite recipe and original table data are supplied in
\texttt{coefficients/upstream/}; \texttt{coefficients/ports.json}
binds the four materialized gzip payloads by full SHA-256 hashes and
their original names. Their uncompressed sizes are
\begin{center}
\begin{tabular}{@{}lrr@{}}
\toprule
Port & Bit length & Bytes\\
\midrule
$\beta$ & $1\,064\,921$ & $133\,116$\\
$K_0$ & $1\,070\,004$ & $133\,751$\\
$M_C$ & $1\,064\,921$ & $133\,116$\\
$M_F$ & $1\,064\,922$ & $133\,116$\\
\midrule
Total & & $533\,099$\\
\bottomrule
\end{tabular}
\end{center}
These integers do not depend on whether the machine halts. They are
computed from its finite transition table and the fixed compiler layout.
The nine- and ten-variable circuits additionally use fixed coefficient
leaves obtained by explicit integer recipes from these four ports.
Every such leaf has a name, recipe, and checked bound. None is an
unspecified index supplied by an existence theorem.

We deliberately retain this common tuple to isolate the new arithmetic
reductions. Package 44 supplies smaller native coefficients, but we do
not claim a coefficient-size record or silently rebind an existing
compression constant to that different tuple.

\subsection{The semantic chain being imported}

The inherited chain identifies three statements: a finite ZFC refutation
exists; the literal blank K215 machine halts; and the calibrated positive
compiler source has a solution. The machine endpoint is represented in
the repository by the theorem
\begin{center}
\small\texttt{\detokenize{KEndpoint215.blank_halts_iff_zfc_refutation}}.
\end{center}
The ordinary mathematical compiler proofs, their source notices, and
the selected coefficient adapter are preserved in this archive.
The scalar Pell arguments are formalizable in ordinary arithmetic
available inside ZFC; this observation is not an exported object-level
proof term.

For the eleven-variable construction we use the following precise
properties of the calibrated source:
\begin{enumerate}
\item From positive coordinates and the transport equation $N_t=1$,
the weak bounds in Lemma~\ref{lem:weak} hold.
\item The normalized Pell equations, including the reconstructed positive
quotients $h,\gamma$, have signed-quotient soundness: a negative recovered
strong root can be accompanied by the corresponding sign change of its
quotient coordinate, and an accepting blank computation is still extracted.
\item A finite blank halt supplies canonical positive normalized
witnesses with main Pell index $R$ and first Pell index $(R+1)/2$.
\end{enumerate}
These are the interfaces used by packages 38 and 44. We do not infer
them from finite modular samples. For low-variable encoding, the two
literal 22-equation quadratic parents are audited separately: package
40's parent and package 45's parent have the same monomial support but
one different row after their affine coordinate choices. The archive
checks the coordinate relation and does not replace one parent by the
other without that transformation.

\section{Support-aware quadratic packing}
The main low-variable improvement is to replace full Sidon separation by
separation only where the input system actually has a nonzero monomial.
This changes the largest fixed exponent label from $799$ to $190$.

\subsection{The exact hypothesis needed}
Let $p_1,\ldots,p_m\in\mathbb Z[X_1,\ldots,X_n]$ have total degree at most
two, let $\mathcal A$ be the union of their supports, and set
$\mathcal Q=\{\alpha\in\mathbb N^n:|\alpha|\leq2\}$.
For distinct positive integers $w_i$, write
$e(\alpha)=\sum_i\alpha_iw_i$ and $W=\max_iw_i$.
The sufficient condition is
\begin{equation}\label{eq:support-isolation}
 \alpha\in\mathcal A,\quad\beta\in\mathcal Q,\quad
 e(\alpha)=e(\beta)\quad\Longrightarrow\quad\alpha=\beta.
\end{equation}
The map $e$ need not be injective on all of $\mathcal Q$.
Let $h_j=2W+(j-1)(2W+1)$, and define
\begin{align*}
 c_{j,\alpha}&=(2-|\alpha|)!\prod_i\alpha_i!\,[X^\alpha]p_j,\\
 D(T)&=\sum_{j,\alpha}c_{j,\alpha}T^{h_j-e(\alpha)},&
 C_z(T)&=\left(1+\sum_i z_iT^{w_i}\right)^2.
\end{align*}
Then all exponents of $D$ are nonnegative and
\begin{equation}\label{eq:support-extraction}
 [T^{h_j}]C_z(T)D(T)=2p_j(z)
 \qquad(1\leq j\leq m).
\end{equation}

\paragraph{Proof.}
A source monomial $\alpha$ of row $k$ and an expansion monomial $\beta$
can contribute to $h_j$ only if
$h_j-h_k=e(\beta)-e(\alpha)$.
The right side has absolute value at most $2W$; a nonzero left side has
absolute value at least $2W+1$. Hence $k=j$, and
\eqref{eq:support-isolation} gives $\beta=\alpha$.
The factorial in $c_{j,\alpha}$ cancels the multinomial denominator,
giving \eqref{eq:support-extraction}. The row intervals
$[h_j-2W,h_j]$ are disjoint and the supported exponents within each are
distinct. Therefore, if $|c_{j,\alpha}|<M$, each coefficient of $D$ has
absolute value less than $M$. Each coefficient of $C_zD$ has absolute
value at most
\begin{equation}\label{eq:support-bound}
 M\left(1+\sum_i|z_i|\right)^2.
\end{equation}
Indeed, the sum of the nonnegative coefficients of $C_{|z|}$ is exactly
the squared quantity. Collisions between absent monomials preserve this
sum. In particular, no additional factor $m$ is needed.
\hfill$\square$

\subsection{The explicit instance}
The variable order is the inherited order
\[
 J,F,\alpha,z,f,h,i,o,j,s,w,\tau,\eta,\zeta,y,\gamma,
 N,E,P,c,a,X,D,C_2,L,F_2,V,g,I,B_{\rm aux},\Delta.
\]
In that order the weights are
\begin{align*}
 &(10,14,128,112,178,157,173,145,129,63,123,84,151,75,2,124,\\
 &\hspace{9mm}3,169,31,190,185,184,189,177,90,111,19,30,68,51,144).
\end{align*}
With a padding weight zero, there are $\binom{33}{2}=528$ degree-two
multisets. The source uses only $71$ of their monomials. Each supported
sum is unique among all $528$ sums, while the entire list has only $278$
distinct values. Thus $250$ repetitions are allowed, and the vector is
explicitly not a $B_2$ set. The finite verifier reconstructs the source
supports from the literal parent circuits and checks all
$136\cdot528=71{,}808$ source/expansion pairs. Its only $136$ target hits
are the required diagonal hits.

Two inherited parent presentations are retained exactly. The ten-integer
construction uses package~40's 65-operation parent; the nine-variable
construction uses package~45's 60-operation parent. Their supports agree,
but their transport coordinates satisfy $z_{\rm old}=z_{\rm new}+1$.
The affine polynomial identity is checked independently. The inherited
positive-zero correspondence, including the step $z_{\rm old}\geq2$, is
recorded in the bundled parent proof. Their coefficients are not silently
identified. Both use the same complete four-port tuple of package~31,
including the already shifted fourth port.

\subsection{Why the unbounded source interface is preserved}
First rewrite the inherited positive-variable parent by $x_i=z_i+1$, so
that the source coordinates $z_i$ are nonnegative. Move its fifth shifted
residual to the last row. Its value at the zero vector is $1$, so the highest coefficient of $D$ is $2$, and the source
has no common zero at the zero vector. Positivity of the weights and a
geometric-series bound give $D(B)>0$ whenever $B>M+1$.
All no-carry arguments use only this positive leading term,
\eqref{eq:support-extraction}, \eqref{eq:support-bound}, and the maximum
weight. They do not need the full Sidon hypothesis.

For the nine-variable form use $h_j=381j$, so
$h_{22}=8382$ and $K=h_{22}+2W=8762$.
The unchanged bound on weighted coefficients is
$M=2^{2134927}$; put $b_0=4096M=2^{2134939}$,
$b=b_0+\varepsilon$ and $B=b^3$. Define
\begin{align*}
 \ell&=\sum_iB^{w_i},& N_0&=B^{191},&N_1&=4B^{8763},&N&=N_0N_1,\\
 M_0&=N_0-1-(b-1)\ell,& T_0&=\sum_jB^{h_j},\\
 T&=M_0+(B-2)T_0N_0,\\
 J_0&=2(1+g)^2D(B)+B\sum_{k=0}^{8762}B^k,& S&=g+J_0N_0,\\
 R&=(N-1)\bigl((S+T+1)N+T+1\bigr),&G_{\rm cap}&=32bB^{190}.
\end{align*}
The bounded branch $g<4bB^{190}$ gives $g<N_0$ and $0<J_0<N_1$;
also $0<T<N$. When $b$ is a power of two, the first mask forces
$g=\sum_i z_iB^{w_i}$ with $0\leq z_i<b$.
Equation~\eqref{eq:support-bound} gives strict coefficient magnitudes
below $M(32b)^2\leq B/4$. After adding $B/2$ to each coefficient, all
digits lie strictly between zero and $B$. The second mask forces every
target digit to equal $B/2$, and therefore every source row vanishes.
The binary-popcount identity gives the inherited central-binomial
divisibility condition $N^2\mid\binom{2R}{R}$.

Conversely, choose an arbitrarily large power of two $b$ above the
coordinates of a common source zero. The resulting packed $g$ is positive,
all target tests pass, and $32g<G_{\rm cap}$. Thus both directions of the
Jones ratio interface hold. For Sun's ten-integer interface use
$h_j=380+381(j-1)$, $K=8761$, $b=1+3f$ and
$B=2^{2140030}b^2$. The original coefficient majorant and all congruence
conditions remain valid. The source cap is still $bB^{190}$, in precisely
the required fixed-power form. The inherited finalizers therefore apply
without increasing the number of witnesses.

\subsection{Resulting fixed polynomials and certified metrics}
The literal ten-integer circuit has $653$ live binary arithmetic operations
and exact total degree $56{,}593{,}864$. The literal nine-positive-variable
circuit has $651$ operations and exact degree $30{,}947{,}584$; its
nonnegative-variable translation has $654$ operations. For the latter,
$t=\deg N=26{,}862$, and the inherited cancellation-resolved identity gives
\[
 \deg F_9=1152t+12W+280=30{,}947{,}584,
\]
with diagonal leading coefficient $2^{2050}$. A separate certificate circuit
has $666$ operations but represents exactly the same polynomial; the
$651$-operation circuit is used for the operation comparison.
All fixed powers in the delivered circuits are expanded into multiplication
gates. Exact degree lower bounds follow from nonzero leading coefficients
at the proved degree upper bound. Modular full-formula comparisons are
reported separately as implementation evidence.

Against the reviewed wave-6 constructions, the degrees fall from
$237{,}520{,}456$ and $129{,}875{,}980$, respectively. The operation savings
are $129$ and $137$. The consistency equivalences retain their imported
compiler/K215, Jones, and (for the integer form) Sun dependencies. This
contribution proves and certifies the new source transformation; it does
not claim global minimality or a new end-to-end proof-assistant certificate.


\section{An explicit equation in eleven positive variables}
\label{sec:eleven}

The eleven supplied variables, and no others, are
\[
 J,F,\alpha,z,i,T,s,w,\eta,\zeta,t\in\Zpos.
\]
Use the four fixed ports above and define the polynomial abbreviations
\begin{align*}
 q&=\beta J+1,& C&=q-F-\alpha,& X&=wq,& Y&=sq^3,\\
 E&=XY,& k&=\eta+\zeta,& c&=kY+\eta,& a&=E+Y,\\
 H&=4a+3,&\Delta&=a^2+H=(a+2)^2-1,
 &\ell&=EkY,& A&=1+\ell(\ell+k),\\
 G&=1+\Delta c^2,& b_0&=X+ac,& I&=ic^2,&D&=1+\Delta I^2,\\
 Q&=\Delta(D-1)=(\Delta I)^2,&d&=Q-1,&v&=cT,&b&=c+RD,
\end{align*}
where the rank and transport polynomials are
\begin{align}
 R&=((q-1)(q-F)+\alpha)(q^2-1)+(M_C+qM_F)J,\label{eq:rank}\\
 N_t&=(K_0+w)C+q-F-z(q-1).\label{eq:transport}
\end{align}
Put
\begin{align}
 U&=4q^4,&r&=k-R-1,& K&=IT,\nonumber\\
 S&=UH,& L&=U+SK,\nonumber\\
 W&=U\bigl[b_0+H(t+b+\ell)\bigr]+H\bigl(2q^4k-r\bigr),\label{eq:carrier}\\
 e&=Q(Dv^2+b^2)-1,&B_{\!a}&=2Qvb.\nonumber
\end{align}
The symbol $K$ in \eqref{eq:carrier} is a computed weight, distinct
from the fixed coefficient $K_0$.

Define
\begin{equation}
 \mathcal N=
 \prod_{\substack{g=\pm\sqrt G,\ \tau=\pm\sqrt A\\ f=\pm\sqrt D}}
 \left[d(W-Lg-S\tau)^2-S^2e+S^2B_{\!a}f\right].
 \label{eq:formalnorm}
\end{equation}
This is an integer polynomial: the product is invariant under each
formal sign change, and only even powers of the three formal roots
remain. This definition does not assume that the square classes are
independent or that there exist eight independent field automorphisms.
The radical symbols are notation for a finite polynomial identity.

Our main equation is
\begin{equation}
 \boxed{\quad P_{11}=N_t(1+\mathcal N)-1=0.\quad}
 \label{eq:p11}
\end{equation}
Every named quantity in \eqref{eq:rank}--\eqref{eq:p11} is computed
from the eleven supplied variables. There is no variable exponent,
divisibility predicate, square root operation, or extra existential
quantifier in the final literal circuit.

\subsection{A division-free scalar formula}
\label{sec:scalar}

The following evaluation specifies \eqref{eq:formalnorm} using only
integer additions, subtractions, and multiplications:
\begin{align*}
 n_d&=dS^2A,&P_H&=S^2QDv^2,&U_H&=S^2Qb^2,\\
 e_H&=P_H+U_H-S^2,& j&=4P_HU_H,\\
 \mathcal T&=2d(W^2+L^2G),& \kappa&=[d(W^2-L^2G)]^2,\\
 \mathcal A&=2(n_d+e_H),&\mathcal B&=(n_d-e_H)^2-j,\\
 a_r&=\mathcal T-\mathcal A,&b_r&=\mathcal B-\kappa,\\
 v_r&=16n_dj-a_r(a_r\mathcal T+2b_r),
 &u_r&=b_r^2-a_r^2\kappa,\\
 \mathcal N&=\kappa v_r^2-u_r(\mathcal T v_r-u_r).
\end{align*}
For a proof, after eliminating the signs of $\tau$ and $f$, the
polynomial in $x=d(W-Lg)^2$ is
\[
 (x^2-\mathcal A x+\mathcal B)^2-16n_dj x.
\]
The two values obtained from the signs of $g$ have sum $\mathcal T$
and product $\kappa$. Reduction modulo
$x^2-\mathcal T x+\kappa$ gives $u_r-v_rx$; its quadratic norm
is the displayed expression. This proves the scalar identity over an
integer polynomial ring, including all square and zero specializations.

\begin{lemma}[The unit guard]
\label{lem:guard}
For every integer substitution, $\mathcal N$ is congruent to a square
modulo four. Hence a root of \eqref{eq:p11} satisfies
$N_t=1$ and $\mathcal N=0$.
\end{lemma}
\begin{proof}
Set $x_0=d(W^2+L^2G)$. Then
$\mathcal T=2x_0$ and
$\kappa=x_0^2-4d^2W^2L^2G$. Thus
\[
 \mathcal N\equiv(u_r-x_0v_r)^2\pmod4.
\]
At a product root, $N_t$ and $1+\mathcal N$ are integer units.
The possibility $1+\mathcal N=-1$ would require
$\mathcal N=-2$, impossible modulo four. Both units are therefore $+1$.
\end{proof}

\section{Root recovery without assuming the missing quotient}
\label{sec:root-recovery}

\begin{lemma}[Weak bounds]\label{lem:weak}
At a positive root of \eqref{eq:p11}, the authentic coefficient
interface and \eqref{eq:transport} give
\begin{gather*}
 q\ge16,\quad C>0,\quad 3q<R<q^4-2,\quad
 X\ge q,\quad Y\ge q^3,\quad E\ge q^4>R+2,\\
 k\ge2,\quad kY<c<k(Y+1),\quad I\ge c^2,\quad a+2<2E.
\end{gather*}
In particular $c>8192>96$. No inequality $k>E$ is used here.
\end{lemma}
\begin{proof}
The transport and rank bounds are the imported weak interface. In
particular, $C\le0$ would make
$N_t\le1-F-(z-1)(q-1)\le0$, contradicting $N_t=1$.
All remaining assertions follow from the positive coordinate definitions;
$a+2=E+Y+2<2E$ uses $X\ge16$ and $Y\ge4096$.
\end{proof}

In the following estimates, write $g=\sqrt G$, $f=\sqrt D$,
and $\tau=\sqrt A$ for the positive real roots, whether or not they
are integers. The elementary inequalities
\begin{equation}
 1<A<\Delta I=\sqrt Q,\qquad
 g<(a+2)c<2Ec,\qquad f<(a+2)I<2EI
 \label{eq:weak-radicals}
\end{equation}
hold. For the first, $c>kY$ gives
\[
 \Delta c^2-A>E k^2Y(2Y^2-1)-1>0,
\]
and $I\ge c^2$. Notice that these estimates do not use an already
recovered main or first Pell index.

\begin{lemma}[The refined auxiliary gap]\label{lem:gap}
Under the weak bounds,
\[
 d(Ig-cf)>b+cf.
\]
Every coherent complex auxiliary root satisfying
\[
 d y^2=Q(vf'-b)^2-1,
 \qquad (f')^2=D,
\]
therefore satisfies $|y|<b+ITg$.
\end{lemma}
\begin{proof}
The exact relation
\[
 (Ig)^2-(cf)^2=I^2-c^2
\]
and $I\ge c^2$, $c\ge2$, give
\[
 Ig-cf=\frac{I^2-c^2}{Ig+cf}
 >\frac{I^2}{4Ig}>\frac{I}{8Ec}.
\]
Also $d>\tfrac12E^4I^2$. Since $R<E$, $D<4E^2I^2$,
and $c<I^2$, one has $b+cf<6E^3I^2$. Consequently
\[
 d(Ig-cf)>\frac{E^3I^3}{16c}>6E^3I^2>b+cf,
\]
where the middle strict inequality uses $I>96c$.

Set $Z=b+Tcf>1$. Even for a complex $y$, its defining equation
implies $d|y|^2\le QZ^2+1$. The exact difference
\[
 [Z(1+1/d)]^2-\frac{QZ^2+1}{d}
 =\frac{QZ^2-d}{d^2}>0
\]
gives $|y|<Z+Z/d$. As $T\ge1$,
$Z/d\le T(b+cf)/d<T(Ig-cf)$. Thus $|y|<b+ITg$.
\end{proof}

This is the compatibility lemma missing from a direct juxtaposition of
packages 38 and 44: the latter's estimate $Ig-cf>1$ used the old
index quotient. The refined inequality is valid even when $k\le E$.

\subsection{The midpoint and selection of the main root}

For the proof only, put $\delta=\ell+k/2$.
The carrier \eqref{eq:carrier} has integer coefficients even if this
temporary rational abbreviation is not yet an integer. Its inner
argument is
\begin{equation}
 x=\frac{W-Lg'}{S}
   =t+b+\delta-\frac rU+\frac{b_0-g'}H-ITg',
 \qquad (g')^2=G.
 \label{eq:innerarg}
\end{equation}
Since $k\ge2$,
\begin{equation}
 \delta\ge\tau,\qquad
 0\le\delta-\tau
 =\frac{k^2/4-1}{\delta+\tau}
 <\frac{k}{8EY}<Ig-cf.
 \label{eq:midpointgap}
\end{equation}
The last inequality follows from Lemma~\ref{lem:gap}'s stronger bound
$Ig-cf>I/(8Ec)\ge c/(8E)>kY/(8E)$.
We do \emph{not} claim $\delta-\tau<1$; $k$ is unbounded.

Each root of the inner elimination polynomial has a coherent form
$x=\tau'+y$ with $(\tau')^2=A$ and the auxiliary relation in
Lemma~\ref{lem:gap}. Thus $|x|<\tau+b+ITg$.
If $g'=-g$, however, \eqref{eq:innerarg} gives
\[
 x=t+b+\delta-r/U+(b_0+g)/H+ITg>b+\tau+ITg.
\]
Indeed $H<5a$, $b_0>ac$, $g>ac$, and $c>kY$ imply
$(b_0+g)/H>2c/5>k>r/U$. This is a contradiction.
If $G$ were nonsquare, actual conjugation of $\Q(\sqrt G)$ would
carry the vanishing rational-coefficient inner polynomial at the
positive branch to the forbidden negative branch. Therefore $g$ is a
positive integer and $x$ is rational. This argument uses one actual
quadratic conjugation, not fictitious independent signs in a splitting field.

\subsection{Exceptional inner roots and rational Pell saturation}

The inner polynomial is
\begin{align*}
 C_x&=d(x^2+A)-e,\\
 F_*(x)&=\left[C_x^2+DB_{\!a}^2-4d^2x^2A\right]^2
                    -4C_x^2DB_{\!a}^2.
\end{align*}
It has $F_*(x)=0$ for the rational $x$ just recovered. If $C_x=0$,
substitution gives the integer identity
\[
 d^2A^2-deA+Q^2v^2b^2D=0.
\]
Reducing modulo $Q$, where $d\equiv e\equiv-1$, gives
$Q\mid A(A-1)$, impossible since $0<A(A-1)<Q$.
Likewise $F_*(0)=0$ would give $Q\mid(A-1)^2$, also impossible.
Therefore both $x$ and $C_x$ are nonzero. Now define
\begin{align*}
 U_x&=C_x^2+DB_{\!a}^2-4d^2x^2A,\\
 f'&=-\frac{U_x}{2C_xB_{\!a}},&
 \tau'&=\frac{C_x+B_{\!a}f'}{2dx},&y&=x-\tau'.
\end{align*}
All denominators have been justified. Direct substitution gives
$(f')^2=D$, $(\tau')^2=A$, and
\[
 dy^2=Q(vf'-b)^2-1.
\]
The rational square roots $f',\tau'$ of integers are integers.

We also need $y\in\Z$. The inherited rational Pell saturation lemma
states that, for integers $m\ge2$ and $V$, if
\[
 m^2V^2-(m^2-1)y^2=1,\qquad y\in\Q,
\]
then $y\in\Z$. Here $m=\Delta I$ and $V=vf'-b$.
For completeness, write $y=p/q$ in lowest terms. Then $q^2\mid m^2-1$,
and the positive nonsquare $D_0=(m^2-1)/q^2$ has integer-coordinate
Pell solutions $(m,q)$ and $(m|V|,|p|)$. Both are powers of its
fundamental integer-coordinate unit. Their abscissas $\chi_j,\chi_l$
satisfy $\chi_j\mid\chi_l$. Reducing and reflecting $l$ modulo $2j$,
the addition formulas and strict growth show that $l$ is an odd
multiple of $j$. The ordinate formulas then give $q\mid p$;
lowest terms force $q=1$. The cases $V=0$ or $p=0$ cannot satisfy
the original equation. It follows that $y$, and hence $x$, is integral.

\section{A small modulus recovers the deleted Pell quotient}
\label{sec:quotient}

Write $A_0=a+2$ and $P_0=2EY+1$. The first Pell equation has
fundamental integer-coordinate unit $(P_0,2)$, while the main equation
has fundamental unit $(A_0,1)$. Thus
\begin{equation}
 |\tau'|=\chi_{P_0}(n),\quad k=2\psi_{P_0}(n),\qquad
 g=\chi_{A_0}(p),\quad c=\psi_{A_0}(p)
 \label{eq:firstmainindices}
\end{equation}
for positive indices $n,p$. In particular $k$ is even. Only now do
we assert that $\delta$ is integral. From the integer carrier identity
\eqref{eq:innerarg}, it follows that
\begin{equation}
 U\mid Hr.\label{eq:weakdiv}
\end{equation}
Cancelling $H$ at this stage would be invalid.

\subsection{Enough index separation without translating \texorpdfstring{$s$}{s}}

Because $c>1$, we have $p\ge2$. If $n=1$, then $k=2$ and
positivity forces $\eta=\zeta=1$, so $c=2Y+1$. But
$c=\psi_{A_0}(p)\ge2A_0>2Y+1$, a contradiction. Hence
\begin{equation}
 n\ge2,\qquad k\ge4P_0>8EY,\qquad
 c>8EY^2>2R,\qquad c>X.
 \label{eq:recoveredlarge}
\end{equation}
The standard growth estimate also gives $c>2p$.
Thus the necessary index separation comes from the integral Pell
equations themselves. No replacement $s_{\rm old}=s+3q$ is necessary.

The strong equation gives
\[
 |f'|=\chi_{A_0}(m),\qquad I=\psi_{A_0}(m)
\]
for $m\ge1$. Strong divisibility and $c^2\mid I$ first give
$p\mid m$. If $m=pj$, composition yields
\[
 \psi_{A_0}(m)=c\psi_{\chi_{A_0}(p)}(j).
\]
The second factor is congruent modulo $c$ to
$j\chi_{A_0}(p)^{j-1}$. Since $\chi_{A_0}(p)$ is a unit modulo
$c$, we obtain $c\mid j$, hence $pc\mid m$ and $m>2p$.

The auxiliary equation is
\[
 (\Delta I)^2V^2-((\Delta I)^2-1)y^2=1,
 \qquad V=c(Tf'-1)-R(f')^2.
\]
Its Pell index $j_0$ is positive and odd. The usual odd-index
abscissa polynomials $Q_v$ satisfy
\begin{gather*}
 \chi_M(2v+1)=M Q_v(M^2),\\
 Q_v(1-A_0^2)=(-1)^v\psi_{A_0}(2v+1),\qquad
 Q_v(0)=(-1)^v(2v+1).
\end{gather*}
Reducing modulo $|f'|$, using
$Q=\Delta((f')^2-1)\equiv1-A_0^2$ and $V\equiv-c$,
gives $\psi_{A_0}(j_0)\equiv\pm c$. Squaring and doubling gives
\[
 \chi_{A_0}(2j_0)\equiv\chi_{A_0}(2p)
 \pmod{\chi_{A_0}(m)}.
\]
Since $0<2p<m$, signed index reduction gives
$j_0\equiv\pm p\pmod m$. To see why the signed step is valid,
reduce an index modulo $2m$ and reflect it to $[0,m]$. A negative
representative cannot equal $\chi_{A_0}(2p)$ modulo $\chi_{A_0}(m)$
because $2\chi_{A_0}(m-1)<\chi_{A_0}(m)$; strict growth then
forces the representative to be $2p$.

Separately, $Q\equiv0\pmod c$ and $V\equiv-R\pmod c$
give $j_0\equiv\pm R\pmod c$. Since $c\mid m$ and
$0<p,R<c/2$, these congruences force
\begin{equation}
 p=R.\label{eq:pR}
\end{equation}
No step used an index quotient before its reconstruction.

\begin{lemma}[The modulus $4q^4$]\label{lem:modulus}
Under the preceding hypotheses, $2n=R+1$, and both $E$ and $U$
divide $r=k-R-1$.
\end{lemma}
\begin{proof}
The polynomial $H=4sq^3(wq+1)+3$ is odd. Any prime dividing
both $q$ and $H$ must be three. If $3\mid q$, then
$H\equiv3\pmod{27}$, so its 3-adic valuation is exactly one.
Consequently
\[
 \gcd(U,H)\in\{1,3\},\qquad
 \mu=U/\gcd(U,H)\ge4q^4/3>R.
\]
The Pell recurrence gives
$\psi_{P_0}(n)\equiv n\pmod{P_0-1}$, hence
\[
 k\equiv2n\pmod{4EY}.
\]
Moreover $U=4q^4\mid4EY=4ws^2q^7$. Combining this with
\eqref{eq:weakdiv} gives $2n\equiv R+1\pmod\mu$.

Since $P_0>A_0$ and $c>k$, monotonicity of the Pell sequences
in their base and index implies $n<p=R$. Thus
$|2n-R-1|<R<\mu$, and the congruence is an equality.
It follows that $r=k-2n$ is divisible by $4EY$, hence by $E$ and $U$.
\end{proof}

Set
\[
 h=r/E,\qquad \gamma=(g-b_0)/H.
\]
The first is a positive integer because $n\ge2$ and
$\psi_{P_0}(n)>n$. The carrier, $U\mid r$, and even $k$
give $g\equiv b_0\pmod H$, so the second is an integer.
Since $g>(a+1)c$ and $c>X$, it is positive. We have now
reconstructed the normalized source equations
\begin{gather*}
 (\tau')^2-\ell(\ell+k)=1,\qquad
 (b_0+H\gamma)^2-\Delta c^2=1,\\
 (f')^2-\Delta i^2c^4=1,\qquad
 Q[c(Tf'-1)-R(f')^2]^2-(Q-1)y^2=1,\\
 k-hE-R=1,\qquad N_t=1.
\end{gather*}
Take positive absolute values of the first and auxiliary roots. If
$f'<0$, change $f'$ and its quotient coordinate $T$ together,
as allowed by the inherited signed-quotient soundness interface.
This extracts a finite blank-K215 halt and hence a ZFC refutation.

\section{Constructive completeness and exact size of the eleven-variable equation}

\subsection{Constructing the new positive coordinate}

Conversely, take a canonical positive normalized compiler witness.
It has positive $g,\tau,f,y,\gamma$ and first/main indices satisfying
$R=2n-1$. The Pell recurrence gives
$r=k-2n\equiv0\pmod{4EY}$, so $r/U$ is a positive integer;
$k$ is even as well. Define
\begin{equation}
 t=\gamma+ITg+\tau+y-b-\delta+r/U,
 \qquad\delta=\ell+k/2.\label{eq:newt}
\end{equation}
The auxiliary Pell equation gives $y\ge|vf-b|\ge b-vf$.
By \eqref{eq:midpointgap},
\[
 t\ge\gamma+T(Ig-cf)+(\tau-\delta)+r/U
   >\gamma+r/U>0.
\]
Thus \eqref{eq:newt} is an integer positive coordinate. Substituting it
into the carrier makes one coherent factor of \eqref{eq:formalnorm}
vanish, while $N_t=1$ is unchanged. This proves the converse.
Together with soundness, it establishes the eleven-variable case of
Theorem~\ref{thm:main} through the stated inherited calibration.

\subsection{Exact degree, including the cancellations}
\label{sec:degree11}

A naive gate-degree propagation gives 1138 for the product. The actual
degree is 954 because of exact norm cancellations. The following
calculation proves the smaller degree, rather than inferring it from
random evaluations.

Use a subscript zero for a highest homogeneous part, and set
\begin{gather*}
 q_0=\beta J,\quad k_0=\eta+\zeta,\quad
 X_0=wq_0,\quad Y_0=sq_0^3,\quad E_0=X_0Y_0,\quad c_0=k_0Y_0,\\
 I_0=ic_0^2,\quad H_0=4E_0,\quad \ell_0=E_0c_0,\quad
 d_0=E_0^4I_0^2,\\
 R_0=q_0^3(q_0-F),\quad B_0=R_0E_0^2I_0^2,\quad
 V_0=I_0T\ell_0,\quad U_0=4q_0^4.
\end{gather*}
Their relevant degrees are
$\deg H_0=6$, $\deg\ell_0=11$, $\deg d_0=46$,
$\deg B_0=38$, and $\deg V_0=23$.

First omit the factor $U$ and the rational shift $-r/U$ from the
carrier, retaining the rational midpoint $\ell+k/2$. Work in the
Laurent series field obtained by scaling all supplied variables by
one indeterminate; its coefficient field is the rational function
field in those variables and the fixed ports. The principal roots
have leading terms $g_0=\tau_0=\ell_0$ and $f_0=E_0I_0$.
The exact square-difference identity gives
\[
 T(Ig-cf)=\frac{iTc_0}{2E_0}+\text{terms of smaller degree}.
\]
For signs $(s_g,s_\tau,s_f)$, four factors with $s_f\ne s_g$
have degree 119 and leading term
$-s_g4H_0^2d_0B_0V_0$.
For the four factors with $s_f=s_g$, the leading terms are
\begin{center}
\begin{tabular}{@{}ccrl@{}}
\toprule
$s_g$ & $s_\tau$ & Degree & Leading term\\
\midrule
$+$ & $-$ & 107 & $4H_0^2B_0d_0\ell_0$\\
$-$ & $-$ & 107 & $4H_0^2B_0d_0\ell_0$\\
$-$ & $+$ & 101 & $4H_0B_0d_0\ell_0$\\
$+$ & $+$ & 97 & $H_0^2B_0d_0(2t-iTc_0/E_0)$\\
\bottomrule
\end{tabular}
\end{center}
To verify the table, define the full unscaled carrier expressions
\[
 W^{\rm u}=b_0+H(t+b+\ell+k/2),\qquad
 L^{\rm u}=1+HIT,
\]
and put $Z=W^{\rm u}-Hb-Hs_\tau\tau-L^{\rm u}s_gg$.
Expand each factor as
\[
 2dHbZ+dZ^2-H^2b^2-H^2QDv^2+H^2
       +s_f2H^2Qbvf.
\]
The leading degree-119 terms cancel exactly when $s_f=s_g$.
In that case the residual linear expression is controlled successively
by $2\ell$, $(b_0+g)/H\sim2\ell_0/H_0$, or
$t-T(Ig-cf)$; the midpoint error $\delta-\tau$ has degree $-9$.
The quadratic cancellation follows from
\[
 dI^2G-QDc^2=Q(I^2-c^2)-I^2G,
\]
so the remaining quadratic terms are below the listed degrees.
This also supplies a uniform upper bound, not just a leading-term guess.

Multiplying the table gives unscaled norm degree
$4\cdot119+107+107+101+97=888$ and leading form
\[
 2^{14}H_0^{15}d_0^8B_0^8V_0^4\ell_0^3
 \left(2t-\frac{iTc_0}{E_0}\right).
\]
The actual norm is, in the localization at $U$,
\[
 \mathcal N=U^{16}\mathcal N_{\rm unscaled}(t-r/U).
\]
Its degree in the carrier coordinate is at most sixteen, so the
denominators clear. Since $\deg(r/U)=0<\deg t$, the substitution
does not alter the highest homogeneous part. Consequently
\begin{equation}
 \mathcal N_{952}
 =2^{14}U_0^{16}H_0^{15}d_0^8B_0^8V_0^4\ell_0^3
 \left(2t-\frac{iTc_0}{E_0}\right)\ne0.
 \label{eq:lead952}
\end{equation}
The denominators in the displayed expression cancel against its
prefactors. Nonvanishing follows from the independent occurrence of
$t$, and holds for every authentic $\beta>0$. Finally
\[
 (N_t)_2=w(q_0-F-\alpha)-zq_0\ne0,
\]
so \eqref{eq:p11} has exact degree $952+2=954$.

\subsection{Operations and an optional system}

The literal schedule has 69 multiplications and 46 additions or
subtractions, totaling 115. The use of $L=U+SIT$ shares the already
computed $S$, and the midpoint term uses $U/2=2q^4$ directly.
The exact two-equation form $N_t=1$, $\mathcal N=0$ uses 111
operations if the final norm is tested by comparing its two terms.
Applying the inherited positive transport shift
\[
 (K_0+w)C=F+z(q-1),\qquad z_{\rm old}=z+1,
\]
gives the 110-operation system of degrees 2 and 952.
Canonical witnesses have $w=2^R/q>q>F$, so $z_{\rm old}\ge2$;
the shifted coordinate is positive. Arbitrary positive shifted solutions
map back by $z_{\rm old}=z+1$.

The checker computes complete univariate polynomials modulo 10007 and
10009 using the actual materialized coefficients. It obtains nonzero
degree-954 coefficients 8551 and 2035 on the two recorded rays, matching
\eqref{eq:lead952}. These are exact finite-field calculations, not
interpolation. The paper supplies the upper bound; the nonzero coefficient
supplies an independent lower-bound certificate for the literal instance.

\section{A fourteen-coordinate equation of degree 56}
\label{sec:fourteen}

The quotient absorber also applies before radical elimination. Combined
with the inherited ordinary strong norm, it gives fourteen strictly
positive unknowns, exact degree~56, and 79 binary arithmetic operations.
The historical $Q_{56}$ has fifteen positive unknowns, degree~56, and
74 operations. Thus the construction exchanges five operations for one
fewer unknown, while retaining the same degree. This is a separate
tradeoff from the equations in eleven or fewer coordinates.

Keep the common definitions of $q,C,X,Y,E,k,c,a,H,\Delta,R,b_0$,
with $Y=sq^3$ and $b_0=X+ac$. The supplied coordinates are
\[
 J,F,\alpha,z,f,i,T,s,w,\tau,\eta,\zeta,y,\gamma>0.
\]
In addition to $U=4q^4$ and $r=k-R-1$, put
\begin{align*}
 V&=c(Tf-1)-Rf^2, & G&=1+\Delta c^2,\\
 N_0&=\tau^2-(XY^2)(XY^2+1)k^2,
 &I&=ic^2,\\
 A_{\rm aux}&=\Delta(f^2-1),
 &N_a&=A_{\rm aux}(V^2-y^2)+y^2,\\
 N_t&=(K_0+w)C+q-F-z(q-1),
 &D_*&=Ub_0+H(U\gamma-r).
\end{align*}
The single ordinary integer polynomial equation is
\begin{equation}
 \boxed{(D_*^2-U^2G)^2+(I^2-A_{\rm aux})^2+(N_a-1)^2+(N_0N_t-1)^2=0.}
 \label{eq:fourteen}
\end{equation}
Every abbreviation is computed, not another unknown. The same complete
package-31 coefficient tuple supplies all four fixed ports. The coordinate
$i$ now has the ordinary strong meaning: the equation is
$I^2=\Delta(f^2-1)$. One must use the ordinary strong-rank theorem
instead of canceling $\Delta$ informally.

\paragraph{An early lower bound on the index.}
At a zero of \eqref{eq:fourteen}, all four residuals vanish. The
consecutive-product negative-Pell descent excludes $N_0=-1$, so
$N_0N_t=1$ gives $N_0=N_t=1$. The authentic transport interface yields
$q\ge16$, $3q<R<q^4-2$, and $E\ge q^4>R+2$.
The main residual makes $g=D_*/U$ a rational square root of the
integer $G$, hence an integer. It is positive: $H<5a$, $b_0>akY$,
and $UY>5$ give $D_*>Ub_0-Hk>ak(UY-5)>0$.
The carrier implies $U\mid Hr$.

Write $A_*=a+2$ and $P_0=2XY^2+1$. Integer-coordinate Pell
classification supplies positive indices $p,n$ with
\[
 c=\psi_{A_*}(p),\quad g=\chi_{A_*}(p),\quad
 k=2\psi_{P_0}(n).
\]
Because $P_0>A_*$ and $c>k$, we have $p>n$.
The recurrence gives $k\equiv2n\pmod{4XY^2}$, and
$U\mid4XY^2$. Also $\gcd(U,H)\le3$: $H$ is odd, every
common prime divisor of $q$ and $H$ divides~3, and $3\mid q$
implies $H\equiv3\pmod{27}$. Therefore
\[
 \mu=U/\gcd(U,H)\ge4q^4/3>R+1.
\]
The divisibility $U\mid H(k-R-1)$ gives
$2n\equiv R+1\pmod\mu$. Since $n>0$ and $0<R+1<\mu$,
\[
 2n=R+1+j\mu\quad(j\ge0).
\]
Thus $n\ge25$ and $p\ge26$ \emph{before} recovering $p=R$.
The elementary recurrence bound
$c\ge(2A_*-1)^{p-1}\ge A_*^{25}$ gives
\[
 c>A_*\Delta^2,\qquad c>2p,\qquad c>2R.
\]
For the last inequality one may also use $c\ge2A_*>2R$.

\paragraph{Ordinary strong-rank recovery.}
Apply the inherited ordinary strong-rank lemma, with its strict size
hypothesis now proved. From
$c=\psi_{A_*}(p)>A_*\Delta^2$ and
$(ic^2)^2=\Delta(f^2-1)$ it gives a positive index $m$ such that
\[
 f=\chi_{A_*}(m),\quad p\mid m,\quad c\mid m,\quad
 m\ge c>2p,\quad I=\Delta\psi_{A_*}(m),\quad c^2\mid f^2-1.
\]
The proof, reproduced among the inherited sources, handles the possible
nonfundamentality of $A_*$ in the squarefree Pell group. Its large-$c$
condition is essential to this invocation. No assertion that $I/\Delta$
is integral follows merely from the original displayed strong equation.

Set $Q=I^2=A_{\rm aux}$. The auxiliary equation becomes
$QV^2-(Q-1)y^2=1$. Its odd Pell index $j_0$ satisfies
$j_0\equiv\pm p\pmod m$ by reduction modulo $f$ and signed
abscissa reduction; it satisfies $j_0\equiv\pm R\pmod c$ by
reduction modulo $c$. The latter uses the recovered divisibility
$c^2\mid f^2-1$. Since $c\mid m$ and $0<p,R<c/2$, it follows
that $p=R$. These are exactly the ordinary signed auxiliary steps in
the inherited compiler proof, with all required bounds available.

Now $1\le n<p=R$, so
$|2n-R-1|<R<\mu$. The earlier congruence forces $j=0$ and
$2n=R+1$. Hence $r=k-2n$ is divisible by both $E$ and $U$.
Define
\[
 h=r/E,\qquad \gamma_{\rm old}=\gamma-r/U=(g-b_0)/H.
\]
Both are integers. The first is positive since $n\ge25$ and
$\psi_{P_0}(n)>n$. Also $k\ge4P_0>X$ gives $c>X$, while
$g>(a+1)c$ implies $g-b_0>c-X>0$, proving
$\gamma_{\rm old}>0$. All fifteen positive coordinates of the
ordinary-strong blank compiler have been restored, so its calibrated
soundness theorem supplies the ZFC refutation.

\paragraph{Constructive converse.}
A positive ordinary-strong compiler witness satisfies $R=2n-1$.
Thus $r=k-2n$ is a positive multiple of $4XY^2$, hence of $U$.
Set $\gamma=\gamma_{\rm old}+r/U$ and retain the other thirteen
coordinates. Then $D_*=Ug$, and all four residuals in
\eqref{eq:fourteen} vanish. No new condition on $s$ or increase of a
Pell index is needed. Canonical normalized witnesses give ordinary
witnesses by the established substitution
$i_{\rm ordinary}=\Delta i_{\rm normalized}$.
This proves equivalent positive solvability and hence the stated
equivalence of unsolvability with $\operatorname{Con}(\mathrm{ZFC})$.

\paragraph{Exact degree and reproducibility.}
The four residual degrees are $26,22,28,24$. For the first use
\[
 D_*^2-U^2G
 =U^2\bigl((b_0+H\gamma)^2-G\bigr)
  -2UHr(b_0+H\gamma)+H^2r^2.
\]
The three terms have degrees $26$, at most $25$, and at most $20$;
the ordinary main residual in parentheses has exact degree~18.
Only $(N_a-1)^2$ reaches degree~56, with highest homogeneous form
\[
 T^4f^8w^4s^8(\eta+\zeta)^4(\beta J)^{28}.
\]
Its $J^{28}\eta^4T^4f^8w^4s^8$ coefficient is
$\beta^{28}\ne0$. The independent standard-library checker expands
the complete main residual into 2,994 monomials with the coefficient
ports symbolic, proves its degree~26, and checks all 79 live gates
and the final nonzero highest form. This finite certificate verifies
the expression and degree. The unbounded semantic dependencies remain
the inherited ordinary compiler and machine calibration.


\section{Verification, integration, and what remains unformalized}
\label{sec:formal-status}

\subsection{The checks and their proper interpretation}

The archive separates finite implementation checks from mathematical
soundness. Its main checks are:
\begin{itemize}
\item symbolic reconstruction of both quadratic parents and all 136
source coefficient records, with a separate check of their affine
transport-coordinate relation;
\item exhaustive isolation of 71 supported monomials against all 528
quadratic pairs, and all $136\cdot528=71\,808$ source/expansion tests
at the target positions;
\item exact fixed-leaf recipes, coefficient bounds, byte hashes,
operation counts, variable counts, and complete circuit liveness;
\item exact norm identities proved algebraically, with independent
formal-radical multiplication diagnostics including square, zero, and
coincident-radicand specializations;
\item exact degree certificates and tests using the actual fixed tuple;
\item direct accepting and rejecting finite controls for the source
packing, where supplied, and the positive coordinate transformations.
\end{itemize}
None of these finite tests attempts to verify that an unbounded machine
never halts. No search for a ZFC contradiction or for a gigantic zero
was performed. An exhaustive combinatorial isolation certificate is
different from a finite sample of an infinite semantic implication;
the former is precisely the finite hypothesis of the encoding theorem.

\subsection{Formalization status}

At the pinned revision the active research records a Lean-checked
nine-natural-variable existence theorem with coefficients selected through
existential representation, and a checked quartic alternative. It
explicitly leaves executable coefficient export, the explicit $P_{67}$
compiler correctness in Lean, and an object-level ZFC derivation open.
Those facts are preserved in the provenance snapshots.

The present article supplies explicit data and ordinary proofs for the
new adapters. It does not silently turn the existing existence-based
Lean polynomial into the literal circuits here, or claim that local
ring identities prove the compiler endpoint. The new degree and
operation certificates are finite certificates for specific polynomials;
the intended consistency interpretation still composes with the
inherited semantic chain described earlier. The mathematical reviews
in \texttt{review/} were produced as independent AI-assisted reviews
within this work, not as external human peer review.

\subsection{Reproducing and incorporating the package}

Run \texttt{python3 verify.py} from the archive root. Python 3.10 or
later and NumPy are sufficient; the support-encoding and fourteen-variable
checks themselves use the standard library. The verification driver
regenerates the literal circuit files and writes an ordered receipt.
The supplied \texttt{Makefile} rebuilds this article with a standard
\texttt{pdflatex} installation. No network or running ProveIt checkout
is needed for these finite checks.

The intended integration unit is the complete archive directory.
Preserve the coherent coefficient tuple, source ledgers, and licenses.
No existing repository files or remote branches were changed in this
work. A suitable destination is a new descriptively named source
directory under the current DiophantineConsistency report collection;
its number should be assigned at intake instead of guessed here.
The source commit is pinned so that later repository edits do not alter
the comparison silently.

\section{Further research questions}

\begin{enumerate}
\item \textbf{How small can the support-dependent weights be?}
The weight-190 certificate is a feasible construction, not an optimum.
A reproducible SAT or integer-programming search could prove lower
bounds for this exact 71-monomial support. That would distinguish a
search improvement from a structural lower bound.
\item \textbf{Can row targets overlap?}
The present theorem retains separated row blocks. The true obstruction
is a finite set of support/expansion differences. Carefully certified
nonuniform target placement may reduce the maximum target without
requiring a globally stronger pair-sum condition.
\item \textbf{Can the quadratic parent be redesigned for its support?}
Affine coordinate changes and choices of intermediate multiplication
registers change which monomials need isolation. Jointly minimizing
parent support and encoding range may outperform optimizing weights
after the parent has been fixed.
\item \textbf{Can more coefficient or guard structure be retained?}
The present carry estimate uses a safe total coefficient bound.
Row-specific signs, shared coefficient expressions, and smaller
independent buffers may reduce the fixed source base or evaluator,
but each change must preserve the malformed-digit rejection argument.
\item \textbf{Can the new quotient lemma be formalized locally?}
The modulus $4q^4$, its gcd bound, the Pell recurrence congruence,
and the short-index argument are concrete targets for Lean. Their
proof should explicitly order root recovery, evenness of $k$, weak
divisibility, and restoration of $h$.
\item \textbf{Can the full semantic bridge be closed?}
The most valuable formalization work is a data-valued compiler that
connects the actual four ports and literal source circuit to the
existing K215 endpoint. A local simplification certificate alone
does not settle that obligation.
\item \textbf{Can the midpoint carrier remove another coordinate?}
The refined root-gap estimate uses actual coupling between Pell
radicands. A further removal must preserve rational-root saturation,
the nonzero exceptional denominators, and positivity in completeness.
No sub-eleven moderate-degree construction is asserted here.
\end{enumerate}

\clearpage
\appendix
\section{A finite native coefficient recipe}
\label{app:coefficients}

This appendix records the selected tuple without an unspecified
enumeration index. The exact ordered tile list consists of the 943
quadruples in \path{coefficients/upstream/upstream/wang943_tiles.json};
the accompanying generator reconstructs it from the 430 literal K215
actions. Slots are west, north, east, south; start is selector zero,
the impossible tile is selector one, and the remaining tiles are
lexicographically ordered, as in package 31.

Let
\begin{align*}
 \mathcal E_c={}&\{0,\ldots,942\}\cup\{983,\ldots,992\}
 \cup\{1003,\ldots,1032\}\\
 &\cup\{1063,1064,1095,1096,2195,3296\}.
\end{align*}
It has 989 elements. Set $W_j=2^{10+2j}$ for $0\le j<44$,
and copy offsets $o_t=(1023,1013,1003,983)$. If tile $i$ has
slot labels $x_{it}$, define
\[
 c_i=1+\sum_{t=0}^3\sum_{j=0}^9
        \operatorname{bit}_j(x_{it}+512)W_{10t+j}
        +W_{40}+W_{41}+W_{42}+W_{43}.
\]
At exponent $o_t+j$ put $c_e=W_{10t+j}$; at the four listed
anchors put $W_{40},\ldots,W_{43}$. The two dummy coefficients
at 1063 and 1064 are zero. Define
\begin{align*}
 \mu_c&=1022+\sum_{j=0}^{43}W_j,\\
 Z_c&=(2^{119}-1)\sum_{e=8792}^{8798}\rho^e
                 +(2^{27}-1)\rho^{8799},\\
 D_C&=\rho^{30}+\rho^{62}+\rho^{1100}+\rho^{2200}
       +\sum_{e\in\mathcal E_c}c_e(\rho^{4393-e}+\rho^{5495-e}),\\
 D_R&=\rho^{42},\\
 M_C&=B_c-1-\sum_{e\in\mathcal E_c\setminus\{1\}}\rho^e
                       -2\rho^{1063},\\
 M_{F,\mathrm{native}}&=4+\mu_c(\rho^{4393}+\rho^{5495})
       +\sum_{j=0}^9(\rho^{1013+j}+\rho^{1065+j})
       +\rho^{2195}+\rho^{3296}+Z_c.
\end{align*}
Use $\rho=2^{121}$, $B_c=\rho^{8801}$ and the four ports
specified earlier. These are bounded finite sums of fixed integers.
The source proof checks carry bounds, clause decoding, alignment,
mask populations, and the controller for the actual packed rank;
they are not new claims established merely by rewriting the recipe here.

\clearpage
\begin{thebibliography}{9}
\raggedright
\bibitem{repo}
V. Reshetnikov and contributors, \emph{ProveIt}, revision \pin,
23 October 2026. Current collection, packages 31, 38, 40, 44, 45,
and active DiophantineConsistency status and short-equation sources.
\url{https://github.com/VladimirReshetnikov/ProveIt/tree/124c07c0b77a9e20e84676514d8b1ff78536ea3f/Computability/DiophantineConsistency}.
Exact relative paths, hashes, and preserved notices are included in
the archive's provenance files.
\bibitem{jones1982}
J. P. Jones, \emph{Universal Diophantine equation},
Journal of Symbolic Logic \textbf{47} (1982), no. 3, 549--571.
\url{https://doi.org/10.2307/2273588}.
The repository's Jones interfaces and explicit finalizer are preserved
with their own source notices.
\bibitem{sun2021}
Z.-W. Sun, \emph{Further results on Hilbert's Tenth Problem},
Science China Mathematics \textbf{64} (2021), 281--306.
\url{https://arxiv.org/abs/1704.03504v7}.
\bibitem{sun2026}
Z.-W. Sun, \emph{Ten unknowns for Hilbert's tenth problem over the integers},
preprint, arXiv:2609.01594v2, 2 September 2026.
\url{https://arxiv.org/html/2609.01594v2}.
The ten-integer theorem and its number-theoretic finalizer are
attributed to Sun; they are not new results of this article.
\end{thebibliography}
\end{document}
